\documentclass[10pt,a4paper]{amsart}
\usepackage[margin=19mm]{geometry}
\usepackage{amsmath,amssymb,mathtools}
\usepackage[scr=rsfs,cal=euler]{mathalfa}
\usepackage{xfrac}
\usepackage[unicode,hidelinks,bookmarks=false]{hyperref}
\newcommand{\ie}{i.e.\ }
\newcommand{\cf}{cf.\ }
\newcommand{\resp}{respectively\ }
\newcommand{\Subsection}{\subsection}
\providecommand{\nobreakdash}{\nobreak\hbox{-}}
\setlength{\emergencystretch}{2em}
\multlinegap0pt
\usepackage{bm}
\usepackage{bbm}
\usepackage{dsfont}
\usepackage{xspace}
\usepackage{cancel}
\usepackage{mathtools}
\usepackage{cleveref}
\usepackage{amsthm}
\makeatletter
\@ifundefined{lemma}{\newtheorem{lemma}{Lemma}}{}
\makeatother
\newcommand{\N}{\mathbb{N}} 
\newcommand{\set}[1]{\left\{#1\right\}} 
\newcommand{\y}{\mathbf{y}}
\title[Theoretical Challenges in Learning for Branch-and-Cut]{Theoretical Challenges in Learning for Branch-and-Cut}
\author{Hongyu Cheng, Amitabh Basu}
\date{Source: arXiv:2601.23249v1 (CC BY 4.0)}
\begin{document}
\maketitle

\noindent\emph{Source and licence.} Theoretical Challenges in Learning for Branch-and-Cut. Version arXiv:2601.23249v1 (CC BY 4.0), distributed under the Creative Commons Attribution 4.0 International licence, as recorded in the arXiv licence metadata for this exact version and shown on the abstract page https://arxiv.org/abs/2601.23249v1. Full rights notice: licenses/arxiv-2601.23249-CC-BY-4.0.txt.\par\smallskip

\noindent\textit{Editorial scope.} Counted unit: the Lemma whose source label is \texttt{lem:block\_values} in the article (source0.tex lines 1289--1305, source label \texttt{lem:block\_values}), delivered together with the article's own complete original proof of it (source0.tex lines 1307--1397, one \texttt{proof} environment of 91 lines). No mathematics is written, repaired or re-derived: statement and proof are the article's own text line for line. Required context retained uncounted: eq:block\_C1 (lines 1239--1244); eq:block\_C2a (lines 1239--1244); eq:block\_C2b (lines 1239--1244); eq:block\_C2c (lines 1239--1244); lem:triangle\_pack (lines 1257--1267). Every definition, display and label of those lines is retained unchanged. Omitted from this package and not redistributed: 1--1238, 1245--1256, 1268--1288, 1306--1306, 1398--1901. Nothing inside a retained excerpt is omitted or altered: every retained body is delivered exactly as the article's lines are written, so no omission receipt of any kind accompanies this package. Citations: the retained excerpts carry no citation command at all, so no bibliography entry of the article is transcribed and no reference is redistributed. Reference adaptation: none was needed and none was made; every reference inside the delivered unit resolves inside it, so no unresolved reference or citation is produced. Printed slips kept literal: no printed word, symbol, hypothesis or number of the counted statement or of its proof was altered. Layout and class: the article's own class is \texttt{svjour3} with packages including fix-cm, inputenc, fontenc, microtype, graphicx, subcaption, booktabs, array, tabularx, makecell, tikz, hyperref, xurl, natbib; the wrapper is the lane's pinned \texttt{amsart} editorial wrapper at 10pt on A4, which restates the theorem-like environments the retained text needs and adds the article's own macro definitions that the retained text needs (\texttt{\textbackslash N}, \texttt{\textbackslash set}, \texttt{\textbackslash y}), with extra wrapper packages (bm, bbm, dsfont, xspace, cancel, mathtools, cleveref). The delivered numbering is the unit's own. Assets: no retained span contains a figure environment, a graphics-inclusion command, a PDF-page inclusion, a table or a TikZ picture; no image file is copied into this package, the package declares no asset dependency and no image file is redistributed with it. Licence: version arXiv:2601.23249v1 of this article is distributed under the Creative Commons Attribution 4.0 International licence, as recorded in the arXiv licence metadata for this exact version and shown on the abstract page https://arxiv.org/abs/2601.23249v1. A full rights notice is delivered at licenses/arxiv-2601.23249-CC-BY-4.0.txt. Characters: every retained excerpt is pure ASCII, so the delivered engine, XeLaTeX with its default Latin Modern text font, reports no missing character.
	\begin{align}
		2b_i+p_i            & \le 2,\label{eq:block_C1}  \\
		b_i+y_{i,1}+y_{i,2} & \le 2,\label{eq:block_C2a} \\
		b_i+y_{i,1}+y_{i,3} & \le 2,\label{eq:block_C2b} \\
		b_i+y_{i,2}+y_{i,3} & \le 2.\label{eq:block_C2c}
	\end{align}

\begin{lemma}\label{lem:triangle_pack}
	For $t\in[0,2]$, consider the polyhedron
	\[
		Y(t)=\set{\y\in[0,1]^3 : y_1+y_2\le t,\ y_1+y_3\le t,\ y_2+y_3\le t}.
	\]
	Then
	\[
		\max_{\y\in Y(t)}\ 4y_1+8y_2+8y_3=10t,
	\]
	and the unique maximizer is $y_1=y_2=y_3=t/2$.
\end{lemma}

\begin{lemma}\label{lem:block_values}
	Let $n\in\N$, let $M=7n+8$, and consider a single block of $I_n$.
	In each item below, the stated LP value is computed under the specified fixings, with all other variables in the block left free.
	\begin{enumerate}
		\item If no variables in the block are fixed, then the LP relaxation has the unique optimum
		      \[
			      b=\frac12,\qquad p=1,\qquad y_1=y_2=y_3=\frac34,
		      \]
		      with LP value $M+27$.
		\item If $b$ is fixed to $0$, then the block LP value is $M+20$.
		      If $b$ is fixed to $1$, then the block LP value is $34$.
		\item If $y_1$ is fixed to $0$ or to $1$, then the block LP value is $M+24$ in both cases.
		\item If $y_2$ is fixed to $0$, then the block LP value is $M+22$.
		      If $y_2$ is fixed to $1$, then the block LP value is $M+26$.
		      The same conclusions hold when $y_3$ is fixed instead of $y_2$.
	\end{enumerate}
\end{lemma}

\begin{proof}
	(1) Fix $b\in[0,1]$.
	Maximizing the objective subject to $2b+p\le 2$ and $p\in[0,1]$ sets
	$p=\min\set{1,2-2b}$.
	Let $t=2-b\in[1,2]$.
	Constraints \eqref{eq:block_C2a}--\eqref{eq:block_C2c} are equivalent to $\y\in Y(t)$, so
	\cref{lem:triangle_pack} yields
	\[
		\max_{\y\in Y(t)}\ (4y_1+8y_2+8y_3)=10t,
	\]
	attained uniquely at $y_1=y_2=y_3=t/2$.
	Therefore, for fixed $b$ the block LP value equals
	\[
		g(b)=24b+M\min\set{1,2-2b}+10(2-b).
	\]
	If $b\le 1/2$ then $g(b)=M+20+14b$, which is maximized on $[0,1/2]$ at $b=1/2$.
		If $b\ge 1/2$ then $g(b)=2M+20+(14-2M)b$, which is maximized on $[1/2,1]$ at $b=1/2$ since $M\ge 15$.
	Thus $b=1/2$ is the unique maximizer, and the corresponding optimum is $p=1$ and
	$y_1=y_2=y_3=3/4$, yielding LP value $M+27$.

	\smallskip
	(2) If $b=0$, then \eqref{eq:block_C1} permits $p=1$, and \eqref{eq:block_C2a}--\eqref{eq:block_C2c}
	permit $y_1=y_2=y_3=1$. This yields value $M+20$ and is optimal since all objective coefficients are
	nonnegative.
	If $b=1$, then \eqref{eq:block_C1} forces $p=0$ and \eqref{eq:block_C2a}--\eqref{eq:block_C2c} reduce
	to $Y(1)$. By \cref{lem:triangle_pack}, the maximum $y$-value is $10$, so the value is $24+10=34$.

	\smallskip
	(3) We first treat $y_1=0$.
	Constraints \eqref{eq:block_C2a} and \eqref{eq:block_C2b} become $b+y_2\le 2$ and $b+y_3\le 2$, and
	\eqref{eq:block_C2c} becomes $b+y_2+y_3\le 2$.
	Fix $b\in[0,1]$.
	Maximizing over $p$ again gives $p=\min\set{1,2-2b}$.
	The remaining constraints imply $y_2,y_3\in[0,1]$ and $y_2+y_3\le 2-b$, so the maximum of
	$8y_2+8y_3$ equals $8(2-b)$, attained whenever $y_2+y_3=2-b$.
	Hence the block LP value for fixed $b$ equals
	\[
		h_0(b)=24b+M\min\set{1,2-2b}+8(2-b).
	\]
	If $b\le 1/2$ then $h_0(b)=M+16+16b$.
	If $b\ge 1/2$ then $h_0(b)=2M+16+(16-2M)b$.
	In both cases the maximum is attained uniquely at $b=1/2$, yielding $p=1$ and block LP value $M+24$.
	One optimal choice sets $y_2=y_3=3/4$.

	Now treat $y_1=1$.
	Constraints \eqref{eq:block_C2a} and \eqref{eq:block_C2b} become $b+y_2\le 1$ and $b+y_3\le 1$, and
	\eqref{eq:block_C2c} becomes $b+y_2+y_3\le 2$.
	Fix $b\in[0,1]$.
	Maximizing over $p$ gives $p=\min\set{1,2-2b}$.
	The first two inequalities imply $y_2,y_3\le 1-b$, so the maximum of $8y_2+8y_3$ equals $16(1-b)$,
	attained uniquely at $y_2=y_3=1-b$.
	Hence the block LP value for fixed $b$ equals
	\[
		h_1(b)=24b+M\min\set{1,2-2b}+4+16(1-b).
	\]
	If $b\le 1/2$ then $h_1(b)=M+20+8b$.
	If $b\ge 1/2$ then $h_1(b)=2M+20+(8-2M)b$.
	In both cases the maximum is attained uniquely at $b=1/2$, yielding $p=1$, $y_2=y_3=1/2$, and block LP value $M+24$.

	\smallskip
	(4) We analyze branching on $y_2$; the case with $y_3$ is the same by symmetry.
	Fix $y_2=0$ and assume no other variables in the block are fixed.
	Then \eqref{eq:block_C2a} becomes $b+y_1\le 2$ and \eqref{eq:block_C2c} becomes $b+y_3\le 2$, while
	\eqref{eq:block_C2b} becomes $b+y_1+y_3\le 2$.
	Fix $b\in[0,1]$.
	Maximizing over $p$ gives $p=\min\set{1,2-2b}$.
	The remaining constraints imply $y_1,y_3\in[0,1]$ and $y_1+y_3\le 2-b$, so the maximum of
	$4y_1+8y_3$ equals $4(1-b)+8$, attained at $y_3=1$ and $y_1=1-b$.
	Therefore the block LP value for fixed $b$ equals
	\[
		24b+M\min\set{1,2-2b}+4(1-b)+8.
	\]
	If $b\le 1/2$ then this equals $M+12+20b$, which is maximized on $[0,1/2]$ at $b=1/2$.
		If $b\ge 1/2$ then this equals $2M+12+(20-2M)b$, which is maximized on $[1/2,1]$ at $b=1/2$ since $M\ge 15$.
	Thus the block LP value is $M+22$.

	Next fix $y_2=1$.
	Then \eqref{eq:block_C2a} becomes $b+y_1\le 1$ and \eqref{eq:block_C2c} becomes $b+y_3\le 1$, while
	\eqref{eq:block_C2b} becomes $b+y_1+y_3\le 2$.
	Fix $b\in[0,1]$ and maximize over $p$ to get $p=\min\set{1,2-2b}$.
	The first two inequalities imply $y_1,y_3\le 1-b$, so the maximum of $4y_1+8y_3$ equals $12(1-b)$,
	attained uniquely at $y_1=y_3=1-b$.
	Therefore the block LP value for fixed $b$ equals
	\[
		24b+M\min\set{1,2-2b}+8+12(1-b).
	\]
	If $b\le 1/2$ then this equals $M+20+12b$, which is maximized on $[0,1/2]$ at $b=1/2$.
		If $b\ge 1/2$ then this equals $2M+20+(12-2M)b$, which is maximized on $[1/2,1]$ at $b=1/2$ since $M\ge 15$.
	Thus the block LP value is $M+26$.
\qed
\end{proof}

\end{document}
